% =====================================================================
% Drift-scan feasibility for SKA-Mid-like single-dish intensity mapping
%
% SKELETON. Numbers, equations and figures are in place; the prose is not.
%
% Figures. \graphicspath searches ../../figures/ FIRST -- i.e.
% examples/SKA/figures/, where both plotting scripts write -- then the local
% figures/ as a fallback. So there is no copy step: regenerate and rebuild.
%     /home/geoff/gibbs_venv_312/bin/python ../../ska_hi_plots.py   # 001-006
%     /home/geoff/limTOD/.venv/bin/python   ../../ska_hi_design.py  # design set
% The two venvs are NOT interchangeable: ska_hi_plots.py builds an HI cube and
% needs pyccl, which only gibbs_venv_312 has; ska_hi_design.py only reads
% cached .npz and runs under either.
% The local figures/ holds 2026-09-21 snapshots, now superseded -- it is kept
% only so the paper still builds if a figure stops being generated. Delete it
% once that no longer matters.
% NOTE both directories are gitignored (the root .gitignore has a bare
% `figures/'), so a clean checkout will not build until the figures are
% regenerated.
%
% mnras.cls ALREADY loads geometry and hyperref -- loading either again is the
% "Option clash" error. It also takes no font-size option, so no 11pt.
% graphicx is NOT loaded by default, so load it explicitly -- there is no clash.
% Do NOT use mnras's own `usegraphicx' class option: that branch calls
% \usepackage inside the options section, which modern LaTeX kernels reject
% with "\RequirePackage or \LoadClass in Options Section".
% To go back to a plain article: swap the class line for
%     \documentclass[11pt,a4paper]{article}
% and add \usepackage[margin=2.4cm]{geometry} \usepackage[hidelinks]{hyperref},
% and change every figure* back to figure (mnras is two-column, article is not).
% =====================================================================
\documentclass[a4paper,fleqn,usenatbib]{mnras}

\usepackage{graphicx}
\usepackage{amsmath,amssymb}
\usepackage{booktabs}
\usepackage{siunitx}

% siunitx has no parsec, so \mega\parsec parses as a prefix with no unit --
% that is the "Found prefix part with no unit" error.
\DeclareSIUnit{\parsec}{pc}

\graphicspath{{../../figures/}{figures/}}

\newcommand{\kpar}{k_{\parallel}}
\newcommand{\kperp}{k_{\perp}}
\newcommand{\Tb}{T_{\mathrm{b}}}

% mnras wants a short running title and a short author list in the brackets.
\title[Drift scanning is information-limited]{Drift scanning for SKA-Mid-like
single-dish intensity mapping: an information-limited regime}

\author[G. Murphy et al.]{
G. Murphy$^{1}$\thanks{E-mail: TODO}
and \textit{(co-authors)}
\\
$^{1}$TODO affiliation
}

\begin{document}
\maketitle

\begin{abstract}
% TODO. Skeleton of the argument:
% (i) a parked dish cannot cross-link at any azimuth -- geometric, not a
%     matter of degree;
% (ii) the residual is set by the number of measured sky modes, not by
%      sensitivity, and integration does not change it. QUOTE tr(WA) PER UNIT
%      SKY -- 14.6 against 22.3 per 100 deg^2 -- NOT the old "17 against 34",
%      which is the 1%-of-lambda_max count that Sec. 4 demotes to a robustness
%      check because it ranks a 12-strip drift above the raster (52 vs 33)
%      while that drift's floor is 9.3x worse. The three framings give ratios
%      of 1.53, 3.3 and 2.0; only the first is the one the paper defends.
% (iii) with an injected HI signal, nothing is recoverable at the bottom of
%       Band 1 -- the post-clean residual is 347x (drift) and 24x (raster)
%       the signal, and it is the beam + prior floor, not noise;
% (iv) at the bottom of Band 1 cross-linking improves every part of this by
%      roughly an order of magnitude, which is the practical recommendation --
%      but see (v): this is the one claim that is band-specific.
% (v) repeating at 500-550 and 675-725 MHz keeps the negative result at every
%     band, yet INVERTS the ordering: the drift's post-clean foreground
%     residual falls 105x across Band 1 against the raster's 2.5x, crossing
%     over near 600 MHz.
\end{abstract}

% =====================================================================
\section{Introduction}
% TODO
% Context: single-dish IM, MeerKLASS, SKA-Mid Band 1. Why drift scanning is
% attractive (no moving parts, stable ground/spillover signature) and what
% the open question is.

% =====================================================================
\section{Simulations}
\label{sec:sims}

\subsection{Instrument and sky}
Single \SI{15}{\metre} dish at the Karoo site
($\phi = \ang{-30.713}$), \SI{2}{\second} sampling, \SI{3}{\hour} total
integration, HEALPix maps at $N_{\mathrm{side}} = 64$. The sky is GSM2008
rotated to equatorial coordinates. Two beams: a symmetric Gaussian with
$\mathrm{FWHM} = 1.22\lambda/D = \ang{3.99}$ at \SI{350}{\mega\hertz}, and a
tapered circular aperture ($-\SI{14}{\deci\bel}$ edge taper,
$\mathrm{FWHM} = \ang{3.84}$, first sidelobe at $\sim-\SI{23}{\deci\bel}$).

\subsection{Time-ordered data}
With no receiver temperature and no additive $T_{\mathrm{sys}}$ term, so that
the only signal is sky:
\begin{equation}
  d(t) = \bigl[ A\,s \bigr](t)\,\bigl(1 + g(t)\bigr)\bigl(1 + \eta(t)\bigr),
  \label{eq:tod}
\end{equation}
where $A$ is the pointing-and-beam operator, $g$ is $1/f$ gain drift from a
pure $f^{-2}$ flicker model, and $\eta$ is white noise of fractional variance
$\num{2.5e-6}$. Both noise terms are \emph{fractional}, so the contamination
they inject scales with sky brightness.
% TODO: note the seeding scheme -- identical realisations across scenarios is
% what makes every comparison in this paper controlled.

\subsection{Scan strategies}
Both produce exactly 5400 samples, so comparisons are at fixed clock time.
\textbf{Drift}: parked at azimuth \ang{0}, elevations \ang{52}/\ang{50}/\ang{48}
on three consecutive sidereal days. \textbf{Raster}: MeerKLASS-style
constant-elevation scan at \ang{38.19}, triangle-wave azimuth $\pm\ang{7.5}$
about \ang{45} (rising) and \ang{315} (setting), \SI{6}{\arcminute\per\second}.

\paragraph{A parked dish cannot cross-link.}
A dish at fixed $(\mathrm{az}, \mathrm{el})$ is stationary in the rotating
Earth frame, so its boresight traces a constant-declination circle swept in
$+\mathrm{RA}$. The track position angle is \ang{90} at \emph{every} azimuth
(verified numerically at az $0/45/90/270/315$); azimuth selects which
declination, not the scan direction. The two raster passes cross at \ang{73.9},
which is the one thing the drift cannot do at any pointing.

\begin{figure*}
\centering
\includegraphics[width=\textwidth]{meerklass_scan_tracks.png}
\caption{The two constant-elevation passes --- rising at az \ang{45}, setting
at az \ang{315} --- with the drift patch of Sec.~\ref{sec:sims} overlaid, so
the comparison is on the same sky. The passes cross at \ang{73.9}. Every drift
track in Fig.~\ref{fig:laddergeom} instead runs at position angle \ang{90},
whatever the azimuth, so the drift has no counterpart to this figure: it is a
property of the strategy, not a parameter of it.}
\label{fig:scantracks}
\end{figure*}

\paragraph{Dish count, and why the declination ladder is the drift's only
geometry lever.}
Adding dishes at the \emph{same} pointing adds rows to $d$ but not to $A$: the
extra samples are repeat measurements of sky modes already in the row space, so
they reduce the variance of the modes the strategy measures and leave the
unmeasured subspace untouched. This is the same statement as
``integrating for longer'', and neither moves the beam + prior floor of
Sec.~\ref{sec:mapmaking}, which depends on $A$ alone. The one exception is dishes pointed at \emph{different}
elevations, which by the argument above observe different declinations and so
do contribute new rows.

Because a parked dish at elevation $\theta$ traces the same declination circle
on every sidereal day, $N$ strips observed on $N$ consecutive sidereal days and
$N$ dishes parked at $N$ elevations for a single pass are the same operator $A$:
each strip covers the same local-sidereal-time window either way. The
three-elevation ladder above is therefore already a
three-dish configuration, and a multi-dish array buys the ladder in wall-clock
time rather than in measured modes. Combined with the impossibility of
cross-linking, this leaves the declination ladder --- the number of strips and
their spacing in units of the beam --- as the only geometric freedom a drift
scan has.

\begin{figure*}
\centering
\includegraphics[width=\textwidth]{hi_ladder_geometry.png}
\caption{The declination ladder as sky geometry. Parked due north,
$\delta = \ang{90} + \phi - \mathrm{el}$, so elevation steps are declination
steps one for one and an hour of Earth rotation sweeps each strip through
\ang{15} of right ascension. \emph{Top}: coverage, drawn one FWHM thick.
\emph{Bottom}: the summed beam response across declination. Below one beam the
strips merge and concentrate measurement where a single strip cannot sample; at
one beam and beyond they separate and the ladder buys sky area instead. The
measured mode density is annotated on each.}
\label{fig:laddergeom}
\end{figure*}

% =====================================================================
\section{Map-making}
\label{sec:mapmaking}

The map-maker is a Wiener filter. With $N$ the noise covariance, $S$ the prior
covariance and $\mu$ the prior mean,
\begin{equation}
  m = \bigl(A^{\top} N^{-1} A + S^{-1}\bigr)^{-1}
      \bigl(A^{\top} N^{-1} d + S^{-1}\mu\bigr).
  \label{eq:wiener}
\end{equation}
Equation~\eqref{eq:wiener} is \emph{affine} in $d$: a fixed linear map
$W = (A^{\top}N^{-1}A + S^{-1})^{-1}A^{\top}N^{-1}$ plus a constant prior term.
Two consequences are used throughout. An injected sky component adds linearly,
\begin{equation}
  m(d + A s) - m(d) = W A s,
  \label{eq:linearity}
\end{equation}
and solving a component's own TOD with $\mu = 0$ returns $WAs$ directly, at one
solve rather than two. Equation~\eqref{eq:linearity} is verified numerically to
$5\times10^{-6}$.

\paragraph{The beam + prior floor.}
Solving with \emph{noiseless} data leaves a residual
\begin{equation}
  r_{\mathrm{floor}} = (I - WA)\,s,
  \label{eq:floor}
\end{equation}
i.e.\ sky the strategy never measured, filled in by the prior. This term does
not integrate down, and it is the quantity that limits everything below.

\paragraph{Prior choice.}
The prior enters Eq.~\eqref{eq:wiener} twice, and the two entries have to be
kept apart because they fail in different ways.

The prior \emph{mean} is flat: $\mu = \bar{s}\,\mathbf{1}$, a constant at the
patch mean. No spatial structure is supplied, so every spatial mode in the
recovered map comes from the data. This matters because the obvious
alternative --- a beam-smoothed version of the true sky --- supplies exactly
the structure the reconstruction is being asked to demonstrate, and it flatters
the strategy that measures least: the drift leans on such a prior $1.40\times$
against the raster's $1.05\times$.

The prior \emph{covariance} is diagonal and isotropic,
$S^{-1} = \sigma_{\mathrm{p}}^{-2} I$, with $\sigma_{\mathrm{p}}^{2}$ set to the
variance of the sky over the patch. It carries no angular or spectral
structure; it is a single scalar, and it is the only quantity in the solve that
retains any knowledge of the truth. That scalar is not innocuous, because it
fixes where data stops dominating prior: in the eigenbasis of
$A^{\top}N^{-1}A$ with eigenvalues $\lambda_i$, a mode is data-dominated when
$\lambda_i \gg \sigma_{\mathrm{p}}^{-2}$ and prior-filled otherwise. It
therefore moves the measured mode count of Sec.~\ref{sec:modes} and the floor
of Eq.~\eqref{eq:floor} together, and any statement of the form ``the survey
measures $n$ modes'' is implicitly conditioned on it.
% TODO (results pass): the numbers from the prior-variance sweep,
% 0.01-1000x the patch variance. Headline: loosening the prior raises the mode
% count and worsens the residual, so the two move in OPPOSITE directions --
% modes gained from the prior are not modes gained from the survey.

\begin{figure*}
\centering
\includegraphics[width=\textwidth]{hi_priorvar.png}
\caption{The trade space is not a picture of the prior. $\sigma_{\mathrm{p}}^2$
is swept over \numrange{e-2}{e3} times the patch variance with everything else
held fixed; note that only the prior \emph{amplitude} varies, the mean being
flat throughout. Loosening the prior raises $N_{\mathrm{eff}}$ and worsens the
residual at the same time, so the two move in opposite directions: directions
reclassified as data-dominated by relaxing $\sigma_{\mathrm{p}}$ are not
directions the survey measured.}
\label{fig:priorvar}
\end{figure*}

% =====================================================================
\section{How much sky does a strategy actually measure?}
\label{sec:modes}

Eigen-analysis of $A^{\top}N^{-1}A$ counts the independently constrained sky
modes. Writing its eigenvalues as $\lambda_i$ and the prior of
Sec.~\ref{sec:mapmaking} as $S^{-1} = \sigma_{\mathrm{p}}^{-2}I$, the map is
$m = WAs + (\text{prior term})$ and the resolution matrix $WA$ has trace
\begin{equation}
  N_{\mathrm{eff}} \;=\; \mathrm{tr}(WA)
  \;=\; \sum_i \frac{\lambda_i}{\lambda_i + \sigma_{\mathrm{p}}^{-2}},
  \label{eq:neff}
\end{equation}
the number of degrees of freedom the \emph{data} constrains: each mode counts
between 0 and 1 according to whether data or prior dominates it.
Equation~\eqref{eq:neff} needs no threshold, and it is exactly complementary to
the floor, since $N_{\mathrm{pix}} - N_{\mathrm{eff}}$ counts the prior-filled
directions that Eq.~\eqref{eq:floor} is built from. Counting instead the modes
above a fixed fraction of $\lambda_{\max}$ is reported below as a robustness
check only: the drift-to-raster ratio drifts with where the line is drawn,
which is why it is not the headline.

Two qualifications govern how Eq.~\eqref{eq:neff} may be compared across
strategies. First, it must be normalised by the observed solid angle: a
narrower beam buys modes while shrinking the patch, so an unnormalised count
rewards observing less sky. Second, and less obvious, modes per unit sky is a
common currency only at \emph{fixed} $\sigma_{\mathrm{p}}$. Relaxing the prior
also raises $N_{\mathrm{eff}}$, by reclassifying weakly constrained directions
as data-dominated rather than by measuring them, and those modes do not carry
the information that modes gained from the scan strategy do.

The count is grid-independent (identical at $N_{\mathrm{side}} = 32$,
64, 128, 256):

\begin{table*}
\centering
\begin{tabular}{lcccc}
\toprule
 & $N_{\mathrm{eff}} = \mathrm{tr}(WA)$ & modes $> 1\%\,\lambda_{\max}$
 & solid angle [deg$^2$] & $N_{\mathrm{eff}}$ per 100 deg$^2$ \\
\midrule
Drift, 3 strips (published) & 38.78  & 17 & 265.2 & 14.62 \\
Drift, 8 strips             & 100.05 & 37 & 499.4 & 20.03 \\
Drift, 12 strips            & 147.18 & 52 & 674.8 & \textbf{21.81} \\
Raster                      & 128.09 & 33 & 574.1 & \textbf{22.31} \\
\bottomrule
\end{tabular}
\caption{Measured modes at \SI{350.78}{\mega\hertz}, one channel, each
strategy over its own patch. The last column is the comparable one, and the
last two rows are why: a twelve-strip parked drift reaches the raster's mode
density to within \SI{2.2}{\percent} while its floor stays $9.3\times$ worse
(Table~\ref{tab:residual}), so mode density is necessary but not sufficient.
The threshold count in column two is reported only to show that it does not
survive the comparison --- it ranks the twelve-strip drift \emph{above} the
raster, 52 against 33, in the opposite order to the residual. The drift row is
computed twice by independent paths, the eigen-analysis of this section and the
ladder sweep of Fig.~\ref{fig:ladder}, giving 38.78 and 38.77.}
\label{tab:modes}
\end{table*}

\begin{figure*}
\centering
\includegraphics[width=\textwidth]{hi_ladder.png}
\caption{What the drift's only geometry lever is worth. Mode count alone, one
channel, no TOD. \emph{Left}: density against strip spacing; the published
\ang{52}/\ang{50}/\ang{48} configuration is marked. \emph{Middle}: density
against strip count out to $N = 20$, with the two configurations for which a
full HI experiment exists annotated by their measured floor. \emph{Right}:
modes and sky area separately, normalised to a single parked dish. The marginal
return is constant --- some 12 modes per 44\,deg$^2$ for every strip
added --- so the density curve in the middle panel flattens because area grows
alongside modes, not because the ladder saturates.}
\label{fig:ladder}
\end{figure*}


% TODO: the consequence -- Tsys/sqrt(N_samples) overstates the information by
% ~19x, because it assumes every pixel is independently measured.

% =====================================================================
\section{Foreground reconstruction}
\label{sec:foreground}

% TODO: the 001--005 results. Suggested content:
% - 1/f is not the limiter: the residual decomposes in quadrature to a
%   21.4 K beam+prior floor, 7.3 K white, 4.8 K 1/f. Removing 1/f entirely
%   buys 1.6%.
% - Cross-linking on a flat prior: floor -32.6% per pixel, -46.4% at beam
%   scale.
% - Integration: 3 h -> infinity moves the total only 21.9 -> 15.7 K, because
%   the asymptote IS the floor. Information-limited, not noise-limited.

% =====================================================================
\section{HI signal recovery}
\label{sec:hi}

\subsection{Injected signal and geometry}
A log-normal HI cube (Gaussian density $\rightarrow$ HI bias $\rightarrow$
log-normal $\rightarrow$ linear RSD $\rightarrow$ $\Tb(z)$) is generated at
$z_{\mathrm{eff}} = 2.80$, giving $\Tb = \SI{0.518}{\milli\kelvin}$ and
$b_{\mathrm{HI}} = 1.559$. It is placed on the HEALPix grid by a
three-dimensional lightcone interpolation: each (pixel, channel) sample is
evaluated at its true comoving position $r(z)\,\hat{n}$.

At $z \simeq 2.8$ a \SI{15}{\metre} dish resolves \SI{401}{\mega\parsec}
transverse, so $\kperp$ reaches only \SI{0.016}{\per\mega\parsec} while $\kpar$
runs \numrange{0.012}{0.192}. The accessible three-dimensional $k$-space is a
sliver near the $\kpar$ axis, and all spectra below are quoted against $\kpar$
for that reason.

\subsection{Estimator}
With $w_i$ a Blackman taper over $N_c$ channels of comoving depth $\Delta r$,
$L = N_c \Delta r$, and $\overline{w^2}$ the mean squared taper,
\begin{equation}
  P_{ab}(k) = \frac{\Delta r^{2}}{L\,\overline{w^{2}}}
    \Bigl\langle \operatorname{Re}\bigl[\tilde a(k)\,\tilde b^{*}(k)\bigr]
    \Bigr\rangle_{\mathrm{pix}},
  \qquad k = 2\pi f_{\mathrm{rfft}} .
  \label{eq:pkpar}
\end{equation}
The taper is not cosmetic: the foregrounds are $\sim10^{4}$ times the HI, so
band-edge leakage would otherwise dominate every high-$\kpar$ bin.

\subsection{Foreground cleaning and the transfer function}
Cleaning removes the leading $n$ eigenmodes of the channel--channel covariance
$C = \operatorname{cov}(x - \bar{x})$. With $U$ their eigenvectors,
\begin{equation}
  x_{\mathrm{clean}} = x - \bigl[ U U^{\top}(x - \bar{x}) + \bar{x} \bigr].
  \label{eq:pca}
\end{equation}
This is the one step that is \emph{not} linear in the data, which is why the
signal loss must be measured by injection rather than computed. Following
\citet{Cunnington2023}, with $\mathcal{C}_n$ the clean at $n$ modes and $X_m$ an
independent mock,
\begin{equation}
  T(k) = \frac{P\bigl(\mathcal{C}_n[X + X_m] - \mathcal{C}_n[X],\ X_m\bigr)}
              {P(X_m,\,X_m)} .
  \label{eq:tf}
\end{equation}
The cross-power in the numerator avoids the positive noise bias an auto-power of
the difference would carry. Because $X_m$ is each strategy's own map-made
response $WAs_m$, $T$ is a ratio within one arm and is therefore
depth-independent.

\subsection{Common resolution}
The beam narrows \SI{12.5}{\percent} across \SIrange{350}{400}{\mega\hertz}.
Each channel is reconvolved to the widest beam in the band with a Gaussian
kernel
\begin{equation}
  \theta_{k}(\nu) = \sqrt{\theta_{\max}^{2} - \theta^{2}(\nu)},
  \label{eq:commonres}
\end{equation}
exact for the Gaussian beam used here. This halves the raster's residual and
leaves the drift unchanged.

% ---------------------------------------------------------------------
\subsection{Results}

\begin{figure*}
\centering
\includegraphics[width=\textwidth]{hi_transfer_function.png}
\caption{Fraction of map-made HI surviving the clean, $T(\kpar)$, per number of
modes removed. Band is $\pm1\sigma$ over 20 independent mocks. Loss concentrates
at low $\kpar$, where the modes are smoothest along the line of sight and least
distinguishable from foreground.}
\label{fig:tf}
\end{figure*}

\begin{figure*}
\centering
\includegraphics[width=\textwidth]{hi_ablation.png}
\caption{Noiseless (``floor'') and full data agree to within a few per cent in
power -- the open rings sit on the line -- while noise alone runs $52\times$
(drift) and $10\times$ (raster) lower in amplitude. The blocker is structural,
not a matter of sensitivity.}
\label{fig:ablation}
\end{figure*}

\begin{figure*}
\centering
\includegraphics[width=\textwidth]{hi_cleaning_depth.png}
\caption{The first question asked of any floor-limited result: clean deeper?
\emph{Left}: post-clean residual against PCA depth. \emph{Right}: the same
residual measured against the HI that \emph{survives} the clean, which is the
quantity that decides whether the extra mode was worth removing. Both fall and
the residual falls faster, so the net never turns over within the range tested,
in any of the three bands. There is no optimal cleaning depth to find: the
limit is the beam + prior floor, not the cleaning.}
\label{fig:cleandepth}
\end{figure*}

\begin{figure*}
\centering
\includegraphics[width=\textwidth]{hi_realisations.png}
\caption{How much of a single-realisation number is the realisation. Nine runs
on the same pixels, geometry and noise seeds, varying only the true HI field.
\emph{Left}: per strategy, with mean and $\pm1\sigma$ --- \SI{13.4}{\percent}
(drift), \SI{8.7}{\percent} (raster). \emph{Right}: the \emph{paired}
drift-to-raster ratio, \SI{6.4}{\percent}, tighter because the realisation
partly cancels when both strategies see the same field. This is the same size
as several effects reported here, and it is why ratios between strategies are
quoted rather than absolute residuals.}
\label{fig:realisations}
\end{figure*}

\begin{table*}
\centering
\begin{tabular}{lccc}
\toprule
residual / HI, 4 modes & as run, 277\,px & interior, 119\,px & common res., 119\,px \\
\midrule
Drift  & 570.6 & 322.1 & 347.3 \\
Raster & 65.9  & 48.5  & \textbf{24.0} \\
\midrule
Drift / raster & $8.7\times$ & $6.6\times$ & $\mathbf{14.4\times}$ \\
\bottomrule
\end{tabular}
\caption{Two corrections, kept separable: restricting to pixels more than
\ang{3} from the patch edge removes an edge inflation, and the reconvolution of
Eq.~\eqref{eq:commonres} then halves the raster's residual while leaving the
drift unchanged. The last row is the drift-to-raster ratio and \emph{not} a
cross-linking gain: these two strategies also differ in measured mode density
(\num{14.6} against \num{22.3} per 100\,deg$^2$), so the ratio
confounds the two. Controlling for density gives $9.3\times$.}
\label{tab:residual}
\end{table*}

% =====================================================================
\section{Why a different cleaning basis cannot help}
\label{sec:rank}

Every subspace method -- PCA, ICA, GMCA, NMF, kernel PCA -- removes a rank-$n$
subspace of the channel--channel covariance and they differ only in how that
subspace is chosen. The relevant measurement is therefore how many modes each
component occupies.

\begin{table}[h]
\centering
\begin{tabular}{lcccc}
\toprule
modes for & \SI{90}{\percent} & \SI{99}{\percent} & \SI{99.9}{\percent}
  & $\lambda_{10}/\lambda_{1}$ \\
\midrule
GDSM foreground   & \textbf{1} & \textbf{1} & \textbf{1} & \num{2.6e-17} \\
Floor, drift      & 2 & 6 & 10 & \num{6.3e-4} \\
Floor, raster     & 1 & 2 & 5  & \num{5.7e-5} \\
HI, drift         & 14 & \textbf{23} & 28 & --- \\
HI, raster        & 17 & \textbf{25} & 30 & --- \\
\bottomrule
\end{tabular}
\caption{The foreground is rank one over a \SI{50}{\mega\hertz} block; the HI is
the only component that is not compressible.}
\label{tab:rank}
\end{table}

\begin{figure*}
\centering
\includegraphics[width=\textwidth]{hi_chromatic.png}
\caption{Why no choice of basis helps. \emph{Left}: modes needed to carry a
given fraction of each component's variance --- the foreground is rank one, the
floor needs 2--6, and the HI needs 23--25 of 32, so it is the one component
that is not compressible. \emph{Middle}: the eigenvectors themselves. The
foreground's leading mode is a clean power law; the floor's carry progressively
more structure, which is what ``chromatic'' means concretely. \emph{Right}:
principal angles between the floor's subspace and the HI's. They are not
orthogonal, so projecting out the floor necessarily takes HI with it --- which
is why cleaning deeper cannot win.}
\label{fig:chromatic}
\end{figure*}

% TODO: the argument. One PCA mode removes the foreground entirely (after which
% the drift residual is still 896x the HI), so from mode 2 onward the filter is
% fighting the instrumental floor rather than foregrounds. The floor is itself
% low-rank, so a subspace method CAN capture it -- but its modes are the smooth
% ones, and the HI's low-kpar power sits in the same few. Mean cosine of the
% principal angles between the rank-6 floor and HI subspaces is 0.60 (drift)
% and 0.74 (raster). No choice of basis escapes an overlap.

% =====================================================================
\section{Survey design}
\label{sec:design}

% TODO: the synthesis. This section is the reason the paper is more than a
% null result -- it says which knobs a survey designer can turn and what each
% is worth, including the ones that are worth nothing.

\begin{figure*}
\centering
\includegraphics[width=\textwidth]{hi_design.png}
\caption{The design trade space at \SIrange{350}{400}{\mega\hertz}.
\emph{Left}: the beam + prior floor against measured sky modes per unit sky.
Adding parked dishes moves along a shallow power law (slope $-1.28$ from three
to twelve dishes); crossing the tracks is a step change of $9\times$ at the
\emph{same} mode density. Mode count is therefore not a sufficient statistic:
cross-linking changes the structure of the null space rather than its size,
because every drift track shares position angle \ang{90}. \emph{Right}: each
lever as a factor on the floor. Integration time, system temperature and
co-pointed dish count are $1.00\times$ exactly --- the residual is
floor-dominated, so none of them can move it.}
\label{fig:design}
\end{figure*}

\begin{figure*}
\centering
\includegraphics[width=\textwidth]{hi_frequency_design.png}
\caption{Why residual/HI must not be read as a trend across bands.
\emph{Left}: the post-clean foreground residual, an instrument quantity, which
improves $105\times$ for the drift and $2.5\times$ for the raster across Band 1
--- the drift overtakes the raster near \SI{600}{\mega\hertz}. \emph{Middle}:
the HI amplitude, which dims $4.5\times$ over the same range; that is cosmology,
not instrument. \emph{Right}: their quotient, which inverts the ordering and is
non-monotonic for the raster although both of its factors are monotonic. Quote
the two separately whenever bands are compared.}
\label{fig:freqdesign}
\end{figure*}

% =====================================================================
\section{Discussion}
% TODO. Points worth making:
% - The scope of the negative result: drift scanning at the BOTTOM of Band 1
%   on 3 h. Foreground-to-HI contrast is 3978 at 350 MHz against 862 at 1 GHz,
%   and the beam is 2.6x more chromatic; this is close to a worst case.
% - But "cross-linking is the lever" is a 350 MHz statement, not a Band 1 one.
%   Three bands (350-400, 500-550, 675-725) all stay floor-limited and nothing
%   comes within 16x of the HI -- so the negative result is band-wide -- while
%   the drift/raster ordering inverts, crossing near 600 MHz. Above that the
%   drift's null space has narrowed enough that it wins on foreground residual.
% - Do not read residual/HI as a trend across bands: it is an instrument
%   number over a cosmology number. The HI dims 4.5x across Band 1, and the
%   raster's residual/HI runs 27 -> 77 -> 38 although both of its factors are
%   monotonic.
% - Why more integration and more dishes do not help: at common resolution the
%   noise sits 52x (drift) and 10x (raster) below the floor in amplitude, so
%   driving it to zero improves the residual by only 4% and 6% respectively.
% - What does help: mode count, i.e. geometry. And modelling the floor rather
%   than filtering it -- Eq.~\eqref{eq:floor} is known exactly in simulation.
% - Relation to MeerKLASS-type analyses: their strongest results are
%   cross-correlations, where a residual uncorrelated with the tracer inflates
%   the variance rather than biasing the answer; and survey volume averages a
%   residual down in a way a 277-pixel patch cannot.

% =====================================================================
\section{Conclusions}
% TODO

% =====================================================================
\bibliographystyle{plainnat}
\begin{thebibliography}{9}
\bibitem[Cunnington et al.(2023)]{Cunnington2023}
Cunnington S., et al., 2023, MNRAS % TODO complete
\bibitem[Wang et al.(2021)]{Wang2021}
Wang J., et al., 2021, MNRAS % TODO complete
\end{thebibliography}

\end{document}
